\section*{Oscillatore Regolato}

\noindent\colorbox{orange}{
  \begin{minipage}{1.0\linewidth}
    \subsection*{Richiesta (a)}
    Discutere se, sulla base di quanto osservato nella configurazione ad anello, si possa applicare
    il principio di cortocircuito virtuale agli ingressi di entrambi gli operazionali.
  \end{minipage}
}
\subsection*{Risposta (a)}
\paragraph{Op-amp U1}
Il ramo invertente in retroazione, modulo transienti il dispositivo
operer\'a in \colorbox{blu}{regione di linearit\'a}.
\paragraph{Op-amp U2}
Sia il ramo invertente che \underline{non}-invertente in retroazione. Tuttavia, la corrente in retro-azione
\underline{non} finisce nell'ingresso \underline{non}-invertente (\'e presente un ``corto''), la retro-azione
avviene tutta nel ramo invertente, modulo transienti il dispositivo operer\'a in
\colorbox{blu}{regione di linearit\'a}. \\

\noindent\colorbox{orange}{
  \begin{minipage}{1.0\linewidth}
    \subsection*{Richiesta (b)}
    Ricavare un'espressione analitica della L-trasformata della funzione di trasferimento $A_1(s)$
    del blocco 1 nell'ipotesi di idealit\'a dell'operazionale. Si assuma inoltre per semplicit\'a che i
    valori dei componenti siano legati tra loro dalle seguenti relazioni: $R_2=R$, $C_2=C$, $R_1=R/5$ e
    $C_1=5C$.
  \end{minipage}
}
\subsection*{Risposta (b)}
Lo schema del circuito nel blocco 1 \'e quello di un \textit{amplificatore invertente con op-amp} e la
funzione di trasferimento \'e nota: ``l'inverso del rapporto fra l'impedenze del ramo di retroazione e
ramo d'ingresso''
\begin{equation}
  \label{eq:oscillatore-regolato-tfunc-1}
  \colorbox{blu}{\boxed{H_1(s) = -\frac{Z_{R}||Z_{C}}{Z_{5C} + Z_{R/5}} = -\frac{5sRC}{(1 + sRC)^2}}}
\end{equation}
\noindent\colorbox{orange}{
  \begin{minipage}{1.0\linewidth}
    \subsection*{Richiesta (c)}
    Dall'espressione ottenuta calcolare frequenza $f_\pi$ e modulo del guadagno $|A_1(f_{\pi})|$ del blocco 1
    e verificare la loro compatibilit\'a con le precedenti misure di cui al punto 1.2.
  \end{minipage}
}
\subsection*{Risposta (c)}
Restringendo la funzione di trasferimento \ref{eq:oscillatore-regolato-tfunc-1} all'asse immaginario
\begin{equation*}
  A_1(j\omega) = -\frac{5jRC\omega}{(1 + jRC\omega)^2}
\end{equation*}
lo sfasamento di $\pi$ si ottiene per un valore reale del rapporto (siccome c\'e un segno ``meno'' davanti),
questo accade solo se si annulla la parte reale del denominatore (siccome il numeratore ha solo parte
immaginaria) e di conseguenza $\omega_{\pi} = (RC)^{-1}$, sostituendo $\omega \to 2\pi f$ si ottiene la
frequenza di sfasamento $\pi$
\begin{equation}
  \label{eq:oscillatore-regolato-frequenza-pi}
  \colorbox{blu}{\boxed{f_{\pi} = \frac{1}{2\pi RC}}}
\end{equation}
e il guadagno alla medesima pulsazione
\begin{equation}
  \label{eq:oscillatore-regolato-gain-1}
  \colorbox{blu}{\boxed{G(\omega_{\pi}) = \left|\frac{5j}{(1 + j)^2}\right| = \left|\frac{5j}{2j}\right| = \frac52}}
\end{equation}

\noindent\colorbox{orange}{
  \begin{minipage}{1.0\linewidth}
    \subsection*{Richiesta (d)}
    Assumendo (al momento) che i diodi siano interdetti, determinare $A_2$, guadagno in tensione
    del secondo blocco, in termini di k (frazione della resistenza del trimmer verso terra),
    verificando che $-1 < A_2 < 1$.
  \end{minipage}
}
\subsection*{Risposta (d)}
La tensione all'ingresso \underline{non}-invertente \'e quella di un partitore di tensione su
$V_2$, il cui rapporto di partizione \'e controllabile da un fattore $k \in [0, 1]$
\begin{equation*}
  V_+ = \frac{kR}{(1-k)R + kR}V_2 = kV_2
\end{equation*}
nell'ipotesi in cui entrambi i diodi siano interdetti allora la tensione all'ingresso invertente
\begin{equation*}
  V_- = \frac{R}{R + R}V_2 + \frac{R}{R + R}V_3 = \frac12(V_2 + V_3)
\end{equation*}
nell'ipotesi di op-amp \textit{ideale} e in regime di operazione \textit{lineare} applicando il
principio di corto-circuito virtuale si ricava l'espressione da cui calcolare il guadagno del
blocco 2
\begin{equation*}
  kV_2 = \frac12(V_2 + V_3)
\end{equation*}
per cui
\begin{equation}
  \label{eq:oscillatore-regolato-tfunc-2}
  \colorbox{blu}{\boxed{A_2 = (2k - 1)}}
\end{equation}
e come suggerito dalla richiesta $-1 < A_2 < 1$ (siccome $k \in [0, 1]$). \\

\noindent\colorbox{orange}{
  \begin{minipage}{1.0\linewidth}
    \subsection*{Richiesta (e)}
    Calcolare per quale regolazione del trimmer (i.e. per quale valore di k) sia soddisfatta la
    condizione di Barkhausen per l'anello.
  \end{minipage}
}
\subsection*{Risposta (e)}
La funzione di loop-gain
\begin{equation}
  \label{eq:oscillatore-regolato-loop-gain}
  \colorbox{blu}{\boxed{L(s) \equiv H_1(s)A_2(s) = \frac{(5 - 10k)sRC}{(1 + sRC)^2}}}
\end{equation}
nell'ipotesi in cui la \textit{condizione di Barkhausen} possa essere soddisfatta, la frequenza di
auto-oscillazione sar\'a $f_\pi$ da \ref{eq:oscillatore-regolato-frequenza-pi}, occorre regolare
$k \in [0, 1]$ per ottenere guadagno unitario
\begin{equation*}
  L(s) \to L(j\omega) \implies L(j\omega_{\pi}) = \frac{(5 - 10k)j}{2j} = 1
\end{equation*}
da cui si ottiene l'impostazione del trimmer
\begin{equation}
  \label{eq:oscillatore-regolato-trimmer}
  \colorbox{blu}{\boxed{k = 3/10}}
\end{equation}

\noindent\colorbox{orange}{
  \begin{minipage}{1.0\linewidth}
    \subsection*{Richiesta (f)}
    Stimare la d.d.p. ai capi dei diodi nelle condizioni di misura di cui al punto 3 e giustificare
    qualitativamente il loro utilizzo.
  \end{minipage}
}
\begin{equation*}
  \Delta V = V_2/2 - V_3
\end{equation*}
dalla funzione di trasferimento, $V_3$ dipende da $k$ e $V_2$
\begin{equation*}
  \Delta V = \left(\frac32 - 2k\right)V_2
\end{equation*}
sostituendo $k \to 3/10$ trovato in \ref{eq:oscillatore-regolato-trimmer}
\begin{equation*}
  \colorbox{blu}{\boxed{\Delta V = \frac{9}{10}V_2}}
\end{equation*}
la differenza di potenziale ai capi del diodo dipendono dall'intensit\'a del segnale d'ingresso,
seguendo lo schema di amplificazione del blocco 2, all'aumentare di $V_2$ diminuisce
l'amplificazione, stabilizzando l'auto-oscillazione nella presenza di perturbazioni esterne.

